\section{Domain instantiations}
\label{sec:instantiation}

Table~\ref{tab:adapters} lists both adapters side by side.

\begin{table}[t]
\centering
\caption{Domain adapters used in this study. Everything outside this table
(embedding model, retrieval, prompting protocol, iterative loop, feasibility
guard and output schema) is shared code. The maritime column is new in this
paper.}
\label{tab:adapters}
\footnotesize
\begin{tabular}{p{3.0cm}p{4.0cm}p{4.0cm}}
\toprule
Adapter element & Aviation (ADS-B) & Maritime (AIS) \\
\midrule
State schema $\mathcal{S}$ & Latitude, longitude, ground speed, track, barometric altitude & Latitude, longitude, speed over ground (SOG), course over ground (COG), heading \\
Untrusted after onset & Latitude, longitude, ground speed, track & Latitude, longitude, SOG, COG \\
Trusted signal $a$ & Barometric altitude (pressure sensor) & Gyrocompass true heading \\
Phase taxonomy $\mathcal{P}$ & Takeoff / cruise / landing & Departure / transit / arrival \\
Verbaliser $\tau$ (prompt) & Flight-phase template (\ref{app:templates}) & Voyage-phase template (\ref{app:templates}) \\
Environment $\mathcal{E}$ & Meteomatics: wind aloft, turbulence indicators & Open-Meteo Marine: waves, swell, current, wind, sea-surface temperature \\
Memory $\mathcal{M}$ & FlightRadar24 history; one index per phase & NOAA Marine Cadastre history; one index per phase and corridor \\
Kinematics $\kappa$ & $\Delta t = 120$\,s; $\alpha = 1.2$ & $\Delta t = 300$\,s (resampled); $\alpha = 1.2$ (Section~\ref{sec:maritime-calibration}) \\
Reference-speed rule & Inferred from consecutive accepted reconstructed positions & Held at the last trusted SOG for the whole blackout (Section~\ref{sec:maritime-speedrule}) \\
LLM backbone & gpt-5.1 (frozen) & gpt-4.1 (frozen) \\
Retrieved neighbours $K$ & 5 & 6 \\
\bottomrule
\end{tabular}
\end{table}

\subsection{Aviation adapter}

The aviation adapter is that of the conference version
\citep{felendler2026trajrag} and is summarised here only for completeness.
The trusted signal is barometric altitude, which comes from the air-data
system rather than the GNSS receiver. Phases (takeoff, cruise, landing) are
assigned from altitude trends. The verbaliser describes the flight phase,
time since takeoff, altitude trend, heading and speed trend and the aligned
wind and turbulence indicators. During a blackout, ground speed and heading
are not read from ADS-B; they are inferred from consecutive accepted
reconstructed positions and used as the reference speed of
Eq.~\eqref{eq:guard}.

\subsection{Maritime adapter}
\label{sec:maritime-adapter}

\paragraph{Trusted signal} On commercial vessels, true heading is normally
supplied by a gyrocompass, which is independent of GNSS, whereas SOG and COG
are computed by the GNSS receiver. We therefore treat heading as the
maritime counterpart of barometric altitude, and SOG and COG as untrusted
after onset. Heading was present on 98--99\% of AIS reports in our corpus,
comparable to the availability of barometric altitude in ADS-B. Vessels
whose heading comes from a satellite compass break this independence
assumption. We return to this point in Section~\ref{sec:threats}.

\paragraph{Phases and memory} Voyages are segmented into departure (leaving
the berth and manoeuvring), transit (steady channel or approach passage) and
arrival (slowing and berthing). One index is built per phase and per
corridor, so the adapter's memory for a port is simply a small set of
indices built from that port's history. This segmentation is computed
offline from each completed voyage's full SOG trend, which is the practical
way to build phase-indexed memory from historical data; it is not an
implementation of the online phase-assignment rule $\pi$ of
Section~\ref{sec:framework}, which is restricted to trusted information
available up to the current step. At evaluation time, each benchmark
voyage's phase label is taken from this offline segmentation, so the
maritime instantiation evaluates reconstruction \emph{within} a known phase
rather than online phase identification; we return to this in
Section~\ref{sec:threats}.

\paragraph{Verbaliser} The voyage template mirrors the flight template but
uses maritime vocabulary and drivers. A representative departure state
reads: ``Vessel type: cargo. 15.0 minutes into the departure phase, the
vessel is tracking northeast through the Galveston Bay entrance channel,
heading 048\textdegree{} with course over ground 045\textdegree{} and speed
over ground 8.4 knots, gradually accelerating away from berth. Sea state is
moderate, with 0.6\,m wave height and a light following swell; a moderate
current of 0.9\,km/h runs with the vessel's heading, and a 12\,km/h wind
blows from the southeast. Motion has been steady, with no abrupt heading
change.'' The environmental provider attaches wave height, direction and
period, swell height, ocean-current speed and direction, sea-surface
temperature and 10\,m wind speed and direction from Open-Meteo Marine
\citep{openmeteo2025}. \ref{app:templates} gives both templates
verbatim, together with a fully worked maritime reconstruction step.

\paragraph{Sampling} Unlike ADS-B, AIS reporting is irregular: Class A
transponders change their rate with speed and navigational status, and
messages are lost to channel contention \citep{itu2014ais}. Rather than
letting the guard handle a variable $\Delta t$, the adapter resamples every
voyage onto a uniform 5-minute grid before segmentation. We verified that
all consecutive points within a segmented phase are exactly 300\,s apart.
This keeps the engine identical across domains, since the irregularity is
absorbed by the adapter's preprocessing.

\paragraph{Reference-speed rule}
\label{sec:maritime-speedrule}
The maritime $\mathrm{SPEEDRULE}$ holds the reference speed $v$ of
Eq.~\eqref{eq:guard} at the last \emph{trusted} SOG -- the value recorded
immediately before onset -- for the entire blackout, rather than updating it
from any post-onset field. This is the deployable rule used throughout
Section~\ref{sec:results}: it uses no information that a real blackout would
withhold. An earlier implementation instead updated the reference speed from
the recorded SOG at each reconstructed step; because SOG is GNSS-derived,
that variant used information a genuine blackout would not provide, and we
do not report it as a result in this paper. Section~\ref{sec:r1-result}
reports the corrected, deployable numbers together with the size of the
effect relative to that earlier, non-deployable variant, as a transparency
check rather than as an additional claimed result.
